\documentclass[11pt]{article}

\usepackage[margin=0.75in]{geometry}
\usepackage{enumitem}
\usepackage{array}
\usepackage{tabularx}
\usepackage{booktabs}
\usepackage{xcolor}
\usepackage{hyperref}
\usepackage{parskip}
\usepackage{tcolorbox}
\usepackage{amssymb}

\definecolor{lightgray}{gray}{0.95}

\newtcolorbox{activitybox}[1]{
  colback=lightgray,
  colframe=black,
  title=\textbf{#1},
  boxrule=0.7pt,
  arc=2pt
}

\newcommand{\writearea}[1]{%
\vspace{0.2em}
\fbox{\parbox[t][#1][t]{0.97\textwidth}{\vspace{0.3em}}}
}

\title{\textbf{Problem Discovery}\\
\large In-Class Team Activity (and Presentation Guide)\\
DAEN 459 - Fall 2026}
\author{}
\date{}

\begin{document}

\maketitle

\vspace{-5em}

\begin{center}
\textbf{Goal for today:} Clarify the problem before discussing development, architecture, or delivery.
\end{center}

\begin{activitybox}{Ground Rule: No Solutions Yet}
For most of today's activity, do \textbf{not} discuss dashboards, databases, APIs, cloud platforms, machine learning models, visualizations, or implementation details.

If someone says, ``We could build \ldots'', ask:

\begin{center}
\textbf{What problem would that solve?}
\end{center}
%Write solution ideas in the \textbf{Parking Lot} at the end of this handout and return to them only after the discovery work is complete.
\end{activitybox}

\section*{Outline}

\begin{center}
\begin{tabularx}{\textwidth}{p{0.1\textwidth} X p{0.4\textwidth}}
\toprule
\textbf{Time} & \textbf{Activity} & \textbf{Output} \\
\midrule
5 min & Initial problem statement & One sentence \\
5 min & Five Whys & Root problem \\
5 min & Stakeholder map & Stakeholder perspectives \\
5 min & Evidence vs.\ assumptions & Three critical assumptions \\
5 min & Failure even if the code works & Two discovery risks \\
5 min & Discovery questions & Key unknowns \\
5 min & Rewrite problem statement & Revised statement \\
5 min & Prepare report-out & 90-second team presentation \\
\bottomrule
\end{tabularx}
\end{center}

\section*{1. Initial Problem Statement}

Use this structure:

\begin{quote}
\textbf{[Stakeholder] struggles with [problem] because [cause/context], resulting in [consequence].}
\end{quote}

\textbf{Your current problem statement:}

\writearea{1.1in}

\section*{2. Five Whys}

Start with your problem statement. Ask ``Why does this matter?'' several times until you reach an underlying need rather than a surface request.

\begin{center}
\begin{tabularx}{\textwidth}{p{0.10\textwidth} X}
\toprule
\textbf{Step} & \textbf{Why does this matter?} \\
\midrule
1 & \\[0.35in]
2 & \\[0.35in]
3 & \\[0.35in]
4 & \\[0.35in]
5 & \\[0.35in]
\bottomrule
\end{tabularx}
\end{center}

\textbf{What appears to be the deeper problem?}

\writearea{0.8in}

\section*{3. Stakeholder Map}

Identify at least four stakeholder perspectives. Remember: the user, data owner, decision maker, and person affected by the output may be different people.

\begin{center}
\begin{tabularx}{\textwidth}{p{0.18\textwidth} X X X}
\toprule
\textbf{Stakeholder} & \textbf{What do they want?} & \textbf{What do they know?} & \textbf{What could go wrong for them?} \\
\midrule
 & & & \\[0.55in]
 & & & \\[0.55in]
 & & & \\[0.55in]
 & & & \\[0.55in]
\bottomrule
\end{tabularx}
\end{center}

\section*{4. Evidence vs.\ Assumptions}

Separate what you \textbf{know} from what you currently \textbf{assume}.

\begin{center}
\begin{tabularx}{\textwidth}{X X}
\toprule
\textbf{What we know} & \textbf{What we assume} \\
\midrule
 & \\[0.55in]
 & \\[0.55in]
 & \\[0.55in]
 & \\[0.55in]
\bottomrule
\end{tabularx}
\end{center}

\textbf{Circle or list the three assumptions most likely to invalidate the project:}

\begin{enumerate}[leftmargin=*]
    \item \rule{0.88\textwidth}{0.4pt}
    \item \rule{0.88\textwidth}{0.4pt}
    \item \rule{0.88\textwidth}{0.4pt}
\end{enumerate}

\section*{5. How Could the Project Fail Even if the Code Works?}

Generate at least five non-programming failure modes.

Examples:
\begin{itemize}[leftmargin=*]
    \item The data do not measure what we think they measure.
    \item The data are incomplete or unrepresentative.
    \item The wrong stakeholder defines success.
    \item Users do not trust or act on the output.
    \item The information arrives too late to support a decision.
    \item The output does not address the real problem.
    \item The problem is real but not important enough to justify the project.
\end{itemize}

\begin{enumerate}[leftmargin=*]
    \item \rule{0.88\textwidth}{0.4pt}
    \item \rule{0.88\textwidth}{0.4pt}
    \item \rule{0.88\textwidth}{0.4pt}
    \item \rule{0.88\textwidth}{0.4pt}
    \item \rule{0.88\textwidth}{0.4pt}
\end{enumerate}

\textbf{Our two most important discovery risks are:}

\begin{enumerate}[leftmargin=*]
    \item \rule{0.88\textwidth}{0.4pt}
    \item \rule{0.88\textwidth}{0.4pt}
\end{enumerate}

\section*{6. Discovery Questions}

Use the following prompts to identify what your team still needs to learn.

\subsection*{Problem}
\begin{itemize}[leftmargin=*]
    \item What is happening today?
    \item Why is it a problem?
    \item How often does it occur?
    \item What is the consequence?
    \item What would happen if nobody solved it?
\end{itemize}

\subsection*{People}
\begin{itemize}[leftmargin=*]
    \item Who experiences the problem?
    \item Who makes decisions related to it?
    \item Who owns or provides the relevant data?
    \item Who may be affected by the project's conclusions?
\end{itemize}

\subsection*{Data}
\begin{itemize}[leftmargin=*]
    \item What evidence do we currently have?
    \item What data would we ideally want?
    \item What data actually exist?
    \item What important concepts might not be captured in the data?
\end{itemize}

\subsection*{Decisions}
\begin{itemize}[leftmargin=*]
    \item What decision should become easier or better?
    \item Who makes that decision?
    \item What information do they currently use?
    \item How quickly must the information be available?
\end{itemize}

\subsection*{Unknowns}
\begin{itemize}[leftmargin=*]
    \item What are we assuming?
    \item Which assumption is most dangerous?
    \item What do we need to learn before committing to this project?
    \item What evidence could convince us to change direction?
\end{itemize}

\textbf{Our three most important unanswered discovery questions:}

\begin{enumerate}[leftmargin=*]
    \item \rule{0.88\textwidth}{0.4pt}
    \item \rule{0.88\textwidth}{0.4pt}
    \item \rule{0.88\textwidth}{0.4pt}
\end{enumerate}

\section*{7. Rewrite the Problem Statement}

Use this structure:

\begin{quote}
\textbf{For [stakeholder], [current situation/problem] creates [measurable consequence]. We need to understand [unknown/question] so that [decision or outcome] can be improved.}
\end{quote}

\textbf{Revised problem statement:}

\writearea{1.2in}

\section*{8. Prepare the 90-Second Team Presentation\\--- or at least, be prepared for it}

Your report-out should focus on what you learned about the problem, not on what you plan to build.

Use this four-part structure:

\begin{enumerate}[leftmargin=*]
    \item \textbf{What we thought:} ``We originally thought our project was about \ldots''
    \item \textbf{What we learned:} ``After discovery, we think the real problem is \ldots''
    \item \textbf{What matters most:} ``The most important stakeholder / assumption / risk is \ldots''
    \item \textbf{What we still need to learn:} ``Before building anything, we need to understand \ldots''
\end{enumerate}

\begin{activitybox}{Presentation Constraint}
Do not spend presentation time describing implementation, architecture, software tools, or proposed features.

The goal is to show that your team's understanding of the \textbf{problem} became more precise.
\end{activitybox}

\subsection*{Presentation Notes}

\textbf{1. What we thought:}

\writearea{0.65in}

\textbf{2. What we learned:}

\writearea{0.65in}

\textbf{3. Most important stakeholder, assumption, or risk:}

\writearea{0.65in}

\textbf{4. What we still need to learn:}

\writearea{0.65in}

%\section*{Class Discussion Questions}
%
%After each team presents, classmates may ask one discovery-oriented question. Good questions include:
%
%\begin{itemize}[leftmargin=*]
%    \item What evidence supports that this is the real problem?
%    \item Who else should be considered a stakeholder?
%    \item Which assumption would be most dangerous if it were wrong?
%    \item What data would you need to validate this problem?
%    \item What would cause you to redefine or abandon the project?
%    \item What decision will someone make differently if this project succeeds?
%\end{itemize}

\section*{End-of-Class Deliverable: Discovery Canvas (fingers crossed!!}

Before leaving, your team should have concise answers for all six boxes below.

\begin{center}
\begin{tabularx}{\textwidth}{X X}
\toprule
\textbf{1. Problem} & \textbf{2. Stakeholders} \\
\midrule
\vspace{0.8in} & \vspace{0.8in} \\
\midrule
\textbf{3. Evidence} & \textbf{4. Assumptions} \\
\midrule
\vspace{0.8in} & \vspace{0.8in} \\
\midrule
\textbf{5. Unknowns} & \textbf{6. Revised Problem Statement} \\
\midrule
\vspace{0.8in} & \vspace{0.8in} \\
\bottomrule
\end{tabularx}
\end{center}

\section*{Follow-Up: Discovery Plan}

For the next class or homework, identify:

\begin{itemize}[leftmargin=*]
    \item \textbf{3 questions} your team needs answered,
    \item \textbf{1 or 2 people} your team should talk to,
    \item \textbf{2 datasets or documents} your team should search, learn, or inspect,
    \item \textbf{1 assumption} your team should test,
    \item \textbf{1 reason} your team might substantially redefine or abandon the current project direction.
\end{itemize}

%\section*{Parking Lot: Solution Ideas}
%
%Write solution ideas here without discussing them until discovery is complete.
%
%\writearea{2.0in}

\end{document}
